\section{Introduction}
Throughout the years, a lot of efforts were spent to study and derive the laws that govern the motion of a pendulum. \\
In many classical experiments, an observer measures directly the necessary physical quantities in order to estimate some other ones. \\
In our setting, We used an STM32F4Discovery (from now simply Discovery) board that, together with its sensors and an algorithm, it is able to compute directly and provide to the observer the requested results while the experiment is performed. \\
The Discovery board is an electronic device equipped with an ARM CortexM4 microcontroller and various on-board devices. For this project, we used the LIS3DSH accelerometer that measures accelerations of the board in the three axes, and a serial port that enables communication between the board and a Personal Computer that eventually displays the results. \\

\subsection{Problem statement}
A pendulum it is a well studied physical system in which a weight is suspended from a pivot and it is able to oscillate freely. It is often difficult to find exact solutions in the general case so usually, some assumptions are made in order to make the problem easier to model and solve. \\
In our case we model the setting as a simple gravity pendulum, i.e. We assumed that:
\begin{itemize}
	\item the rod or cord on which the bob swings is massless, inextensible and always remains taut,
	\item the bob is a point mass,
	\item motion occurs only in two dimensions,
	\item the gravitational field is uniform,
	\item the support does not move.
\end{itemize}
It is trivial to map the abstract case with our concrete one, in particular, the Discovery board is our bob and the USB cable is the rod. \\
The aim of the project is to estimate the length of the cable while the board oscillates.
\subsection{Summary of the work}
Our work can be divided in three distinct phases:
\begin{itemize}
	\item we analyzed the datasheet of on-board accelerometer and developed a driver that it is able to interact with it while exposing a high-level interface that simplified interaction,
	\item we derived the mathematical model of the physical system in order to be able to separate between variables that we can measure and the unknown that we should compute,
	\item we put all together by means of an algorithm coupled with some parameter tuning and complementary solutions as a result of various experiments. 
\end{itemize}
\pagebreak

\section{Design and implementation}
\textbf{Maybe the following text should be inserted in "summary of the work"}

We drafted the algorithm to calculate the length of the wire. We decided to implement this algorithm in a raw simulator built by us in Javascript \cite{javascriptPrototype}. Then we developed it in MIOSIX.\\
We will see the mathematical model, how to calculate the data we need ($\theta_{max}$ and period) using raw data obtained by the accelerometer.

\subsection{Mathematical Model}
We did not considered the friction (since it is to complex to calculate), so the period of the gravity pendulum has the following series as period formula: \\
\begin{center}
	$T = 2 \pi \sqrt{\frac{L}{g}} (1 + \frac{1}{16} \theta_0^2 + \frac{11}{3072} \theta_0^4 + ...)$\\
\end{center}
Stopping at the second term of the series we are able to calculate the length of the wire ($L$) using period and $\theta_0$, i.e. $\theta_{max}$ , we obtain:\\
\begin{center}
	$ L = g \cdot \frac{period^2}{4 \cdot \pi^2 \cdot (1 + \frac{\theta_0^2}{16})^2}$
\end{center}\par
We can see that if we had used only one term of the series we would not have needed to calculate the $\theta_{max}$, but of course there is an error which increases for larger amplitude. In our application we cannot afford this error since, as we will see, in fact there are other imprecisions due to measurement errors and approximations, and the experiment is conducted using large amplitude.\\
We calculated the errors according to the values obtained using 3 members of the series in table \prettyref{tab:errors}.  

\begin{table}
	\small
	\begin{center}
		
		\caption{Series errors compared with the value obtained with 3 terms}
		\begin{tabular}{l|lcc|L{2cm}l}
			Angle & Error with 1 term & Error with 2 terms\\
			\cline{1-3}
			$0$ & 0\% & 0\%\\			
			$\frac{\pi}{4}$ & 4\% & 0.1\%\\			
			$\frac{\pi}{2}$ & 14.9\% & 1.8\%\\			
			$\pi$ & 49\% & 17.7\%\\			
		\end{tabular}
		\label{tab:errors}
	\end{center}
\end{table}

As we will see the max angle of the pendulum is of $\frac{\pi}{2}$ for that value the error is just of 1.8\%, an acceptable value.

%say that we can use one term and we don't need theta, is an observation
%say how to calculate theta and  period
%say the error in % using one term
%say int he next chapter that for a lot of reasons the experiment is doen using large amplitude, say why
%say the reasons of the angle choosen

\textbf{Please make sure you explicitely cite the figures.
\prettyref{fig:polimi_logo} shows the logo of Politecnico di Milano.
Please also make sure you provide each figure with an exhaustive caption.}

\lipsum[1-12]

\begin{figure}
\centering\includegraphics[width=5cm]{polimi}
\caption{POLIMI logo, which is used by this document as a placeholder.}
\label{fig:polimi_logo}
\end{figure}

\section{Experimental evaluation}
\lipsum[1]
\subsection{Experimental setup}
\lipsum[1-2]
\subsection{Results}
\textbf{Please make sure you explicitely cite the tables.
\prettyref{tab:a_complex_table} shows a complex table.
Please also make sure you provide each table with an exhaustive caption.
Captions for tables must be placed before not after the tables.}
\lipsum[1-4]


\begin{table}
\small
\begin{center}

\caption{Summary of the test scenarios.
$M_1\xrightarrow{\alpha}M_2$ means that an application mapping
is changed from $M_1$ to $M_2$ after application $\alpha$
has terminated.}
\begin{tabular}{l|lcc|L{2cm}l}
\cline{2-6}
& \multicolumn{3}{c|}{Description of the workload} & \multicolumn{2}{c}{Cores allocation}\\
\cline{2-6}
Name of scenario & Application & Threads & $\frac{Threads}{Cores}$ & HMP & HMP w/policy \\
\hline

\multirow{2}{*}{LITTLE 1} & ferret$\dagger$ & 1 & \multirow{2}{*}{1.00} & $0-3$ & $0$ \\
 & vips & 3 & & $0-3$ & $1-3 \xrightarrow{\dagger} 0-3$ \\
\hline
\multirow{2}{*}{LITTLE 2} & freqmine$\dagger$ & 2 & \multirow{2}{*}{1.25} & $0-3$ & $0-1$ \\
 & blackscholes & 3 &   & $0-3$ & $0-3 \xrightarrow{\dagger} 0-3$ \\
\hline
\multirow{2}{*}{LITTLE 3} & bodytrack$\dagger$ & 3 & \multirow{2}{*}{1.25} & $0-3$ & $0-1$ \\
 & facesim & 2 &  & $0-3$ & $0-3$ \\
\hline
\multirow{2}{*}{LITTLE 4} & facesim & 3 & \multirow{2}{*}{1.50} & $0-3$ & $0-1\xrightarrow{\dagger}0-3$ \\
 & blackscholes$\dagger$ & 3 &  & $0-3$ & $2-3$ \\
\hline
\hline


\multirow{2}{*}{big 1} & vips & 3 & \multirow{2}{*}{1.00} & $4-7$ & $4-5 \xrightarrow{\dagger} 4-7$ \\
 & ferret$\dagger$ & 1 & & $4-7$ & $6-7$ \\
\hline
\multirow{2}{*}{big 2} & freqmine$\dagger$ & 2 & \multirow{2}{*}{1.25} & $4-7$ & $4-5$ \\
 & blackscholes & 3 &   & $4-7$ & $6-7 \xrightarrow{\dagger} 4-7$ \\
\hline
\multirow{2}{*}{big 3} & facesim & 2 & \multirow{2}{*}{1.25} & $4-7$ & $4-5$ \\
 & bodytrack & 3 &  & $4-7$ & $4-7$ \\
\hline
\multirow{2}{*}{big 4} & facesim & 3 & \multirow{2}{*}{1.50} & $4-7$ & $4-5\xrightarrow{\dagger}4-7$ \\
 & blackscholes$\dagger$ & 3 &  & $4-7$ & $4-7$ \\
\hline

\end{tabular}
  \label{tab:a_complex_table}
\end{center}
\end{table}

\section{Conclusions and Future Works}
\lipsum[1]
